\documentclass[12pt, a4paper]{article}

% ============================================================
%  Paquetes
% ============================================================
\usepackage[utf8]{inputenc}
\usepackage[T1]{fontenc}
\usepackage[spanish]{babel}
\usepackage{amsmath, amssymb}
\usepackage[ruled, vlined, linesnumbered, spanish]{algorithm2e}
\usepackage{geometry}
\usepackage{hyperref}

\geometry{margin=2.5cm}

% ============================================================
%  Metadatos del documento
% ============================================================
\title{Recocido Simulado (SA) --- Pseudocódigo}
\author{}
\date{}

% ============================================================
\begin{document}
\maketitle

% ------------------------------------------------------------
\section{Notación}
% ------------------------------------------------------------

\begin{itemize}
    \item $G$: instancia CARP con conjunto de arcos requeridos $R$.
    \item $s$: solución actual --- un conjunto de rutas que cubre todos los arcos
          de $R$.
    \item $s^{*}$: mejor solución factible encontrada hasta el momento.
    \item $f(s) = c(s) + \lambda \cdot \mathrm{viol}(s)$: objetivo penalizado
          ($\lambda > 0$).
    \item $\mathcal{O}_{\mathrm{intra}} = \{\texttt{relocate\_intra},\,
          \texttt{swap\_intra},\, \texttt{2opt\_intra}\}$: operadores
          intra-ruta.
    \item $\mathcal{O}_{\mathrm{inter}} = \{\texttt{relocate\_inter},\,
          \texttt{swap\_inter},\, \texttt{2opt\_star},\,
          \texttt{cross\_exchange},\, \texttt{or\_opt\_2},\,
          \texttt{or\_opt\_3}\}$: operadores inter-ruta.
    \item $\alpha \in (0,1)$: factor de enfriamiento geométrico.
    \item $p_{\mathrm{inter}} \in [0,1]$: P(elegir inter-ruta) cuando la
          solución es factible.
    \item $\alpha_{\mathrm{inter}} = 0.8$: P(elegir inter-ruta) cuando hay
          violación de capacidad.
    \item $n = |R|$: número de tareas requeridas (variable de escala
          \emph{instance-aware}).
    \item $d_{\max} = \max_{i,j}\,\mathrm{dist}(i,j)$: distancia máxima
          del grafo (camino mínimo más largo entre todos los nodos).
    \item $L = n^{2}$: longitud del plateau (iteraciones por nivel de temperatura).
    \item $\mathit{patience}$: niveles de temperatura consecutivos sin mejora del
          mejor global antes de activar el recalentamiento ($= 0$ desactiva).
    \item $r_{\mathrm{reheat}}$: fracción de $T_{\mathrm{init}}$ a la que se
          sube la temperatura en cada recalentamiento.
    \item $M_{\mathrm{reheat}}$: reheats consecutivos sin mejora antes de parar
          anticipadamente ($= 0$ desactiva).
\end{itemize}

\medskip
\noindent
\textbf{Calibración \emph{instance-aware} (valores por defecto si no se
proporcionan):}
\[
T_{\mathrm{init}} = \frac{20 \cdot d_{\max}}{n}, \qquad
T_{\mathrm{min}}  = \frac{20 \cdot d_{\max}}{n^{2}}, \qquad
L = n^{2}.
\]

\medskip
\noindent
\textbf{Sesgo inter/intra (compartido con TS simple, RTS y ABC simple).}
En cada iteración del bucle interior se calcula $v \leftarrow
\mathrm{viol}(s)$ y se define
\[
p_{\mathrm{efec}} \;=\;
\begin{cases}
  \alpha_{\mathrm{inter}} & \text{si } v > 0, \\
  p_{\mathrm{inter}}     & \text{si } v = 0.
\end{cases}
\]
Con probabilidad $p_{\mathrm{efec}}$ se elige el grupo
$\mathcal{O}_{\mathrm{inter}}$; en caso contrario,
$\mathcal{O}_{\mathrm{intra}}$.

% ------------------------------------------------------------
\section{Pseudocódigo}
% ------------------------------------------------------------

\begin{algorithm}[H]
\DontPrintSemicolon
\caption{Recocido Simulado para CARP (con Recalentamiento)}
\label{alg:sa_simple}

\KwIn{Instancia CARP $G$;\; solución inicial $s_0$;\;
      factor $\alpha$;\; penalización $\lambda$;\;
      $p_{\mathrm{inter}}$;\; opcionalmente: $T_{\mathrm{init}}$,
      $T_{\mathrm{min}}$, $L$ (calibrados si se omiten);\;
      $\mathit{patience}$,\; $r_{\mathrm{reheat}}$,\; $M_{\mathrm{reheat}}$}
\KwOut{Mejor solución factible $s^{*}$}

\BlankLine
\tcp{Calibración instance-aware}
$n \leftarrow |R|$;\quad
$d_{\max} \leftarrow \max_{i,j}\,\mathrm{dist}(i,j)$ \;
$T_{\mathrm{init}} \leftarrow 20 \cdot d_{\max} / n$ \;
$T_{\mathrm{min}}  \leftarrow 20 \cdot d_{\max} / n^{2}$ \;
$L \leftarrow n^{2}$ \tcp{iteraciones por nivel de temperatura}

\BlankLine
\tcp{Precómputo UNA VEZ de las particiones intra/inter/fallback}
Construir $\mathcal{O}_{\mathrm{intra}}$, $\mathcal{O}_{\mathrm{inter}}$,
$\mathcal{O}_{\mathrm{fallback}}$ a partir de los operadores activos \;

\BlankLine
$s \leftarrow s_0$;\quad $s^{*} \leftarrow s$ \;
$T \leftarrow T_{\mathrm{init}}$ \;
$\mathit{niveles\_sin\_mejora} \leftarrow 0$ \;
$\mathit{reheats\_sin\_mejora} \leftarrow 0$ \;
$n_{\mathrm{reheats}} \leftarrow 0$ \;

\BlankLine
\While{$T > T_{\mathrm{min}}$}{

    \tcp{Snapshot del mejor al inicio del nivel (para detectar mejora de nivel)}
    $f^{*}_{\mathrm{nivel}} \leftarrow f(s^{*})$ \;

    \BlankLine
    \For{$\mathit{it} \leftarrow 1$ \KwTo $L$}{

        \BlankLine
        \tcp{Selección del grupo de operadores (sesgo inter/intra)}
        $v \leftarrow \mathrm{viol}(s)$ \;
        \lIf{$v > 0$}{$p_{\mathrm{efec}} \leftarrow \alpha_{\mathrm{inter}}$}
        \lElse{$p_{\mathrm{efec}} \leftarrow p_{\mathrm{inter}}$}
        \eIf{$\mathrm{Uniforme}(0,1) < p_{\mathrm{efec}}$}{
            $\mathit{op} \leftarrow \textsc{SelectUniform}(\mathcal{O}_{\mathrm{inter}})$ \;
        }{
            $\mathit{op} \leftarrow \textsc{SelectUniform}(\mathcal{O}_{\mathrm{intra}})$ \;
        }

        \BlankLine
        $s' \leftarrow \textsc{AplicarOperador}(\mathit{op},\, s)$ \;
        $\Delta \leftarrow f(s') - f(s)$ \;

        \BlankLine
        \tcp{Criterio de Metropolis}
        \If{$\Delta \leq 0$ \textbf{ o } $\mathrm{Uniforme}(0,1) < e^{-\Delta / T}$}{
            $s \leftarrow s'$ \;
        }

        \BlankLine
        \tcp{Actualización del mejor global (solo soluciones factibles)}
        \If{$\textsc{Factible}(s)$ \textbf{ y } $f(s) < f(s^{*})$}{
            $s^{*} \leftarrow s$ \;
        }
    }

    \BlankLine
    \tcp{Paso de enfriamiento geométrico}
    $T \leftarrow \alpha \cdot T$ \;

    \BlankLine
    \tcp{Mecanismo de recalentamiento (Reheat)}
    \If{$f(s^{*}) < f^{*}_{\mathrm{nivel}} - \varepsilon$}{
        \tcp{El nivel produjo mejora: reiniciar contadores}
        $\mathit{niveles\_sin\_mejora} \leftarrow 0$ \;
    }{
        $\mathit{niveles\_sin\_mejora} \leftarrow \mathit{niveles\_sin\_mejora} + 1$ \;
    }

    \BlankLine
    \If{$\mathit{patience} > 0$ \textbf{ y }
        $\mathit{niveles\_sin\_mejora} \geq \mathit{patience}$}{

        \tcp{Guardar el mejor antes del reheat para comparación posterior}
        $f^{*}_{\mathrm{antes\_reheat}} \leftarrow f(s^{*})$ \;

        $T \leftarrow r_{\mathrm{reheat}} \cdot T_{\mathrm{init}}$ \;
        $\mathit{niveles\_sin\_mejora} \leftarrow 0$ \;
        $n_{\mathrm{reheats}} \leftarrow n_{\mathrm{reheats}} + 1$ \;

        \tcp{Criterio de parada por reheats consecutivos sin mejora}
        \eIf{$f(s^{*}) \geq f^{*}_{\mathrm{antes\_reheat}} - \varepsilon$}{
            $\mathit{reheats\_sin\_mejora} \leftarrow
                \mathit{reheats\_sin\_mejora} + 1$
        }{
            $\mathit{reheats\_sin\_mejora} \leftarrow 0$
        }

        \If{$M_{\mathrm{reheat}} > 0$ \textbf{ y }
            $\mathit{reheats\_sin\_mejora} \geq M_{\mathrm{reheat}}$}{
            \textbf{break} \tcp{parada anticipada por reheats estériles}
        }
    }
}

\BlankLine
\Return $s^{*}$
\end{algorithm}

% ------------------------------------------------------------
\section{Notas de implementación}
% ------------------------------------------------------------

\paragraph{Calibración instance-aware.}
Las fórmulas $T_{\mathrm{init}} = 20 d_{\max}/n$ y $T_{\mathrm{min}} = 20
d_{\max}/n^{2}$ escalan la ventana de temperatura con el tamaño de la instancia:
instancias pequeñas (pocos arcos, distancias cortas) reciben ventanas pequeñas;
instancias grandes, ventanas grandes. El plateau $L = n^{2}$ garantiza que cada
nivel ejecute suficientes perturbaciones para evaluar estadísticamente el
comportamiento a esa temperatura.

\paragraph{Reheat.}
El recalentamiento resuelve el estancamiento en mínimos locales cuando $T$ ya es
demasiado baja para aceptar empeoramientos. En términos metalúrgicos: si el metal
queda atrapado en una imperfección cristalina, se vuelve a calentar para dar a
los átomos nueva movilidad. En el algoritmo, se sube $T$ a $r_{\mathrm{reheat}}
\cdot T_{\mathrm{init}}$ y se reinicia el contador de niveles sin mejora. La
mejor solución nunca se descarta. Con $\mathit{patience} = 0$ el reheat está
desactivado y el SA es el clásico de Kirkpatrick (1983).

\paragraph{Parada anticipada por reheats.}
Con $M_{\mathrm{reheat}} > 0$ se puede interrumpir la corrida si $M_{\mathrm{reheat}}$
reheats consecutivos no producen mejora del mejor global. Este criterio complementa
la condición $T < T_{\mathrm{min}}$ y evita gastar tiempo cuando el algoritmo ya
convergió.

\paragraph{Actualización del mejor global.}
Se mantiene el mejor global $s^{*}$ sobre las soluciones \emph{factibles}
únicamente (aquellas con $\mathrm{viol}(s) = 0$). La solución actual $s$ puede
ser infactible (aceptada por el criterio de Metropolis), pero no actualiza $s^{*}$.
Esto garantiza que la salida del algoritmo siempre sea una solución CARP válida.

\paragraph{Precómputo de operadores.}
Las listas $\mathcal{O}_{\mathrm{intra}}$, $\mathcal{O}_{\mathrm{inter}}$ y
$\mathcal{O}_{\mathrm{fallback}}$ se construyen \emph{una sola vez} antes del
bucle exterior y se pasan al helper en cada iteración sin reconstruirlas. Esta
microoptimización evita la creación de listas Python en el camino caliente del
bucle interior (que ejecuta $L \approx n^{2}$ iteraciones por nivel).

\end{document}
